\section{Component Ablations and a Frozen Transfer Protocol}
\label{sec:ablations}
We replay the same native validation observations through the full adaptive
recorder and three single-component removals at each of the four caps (16 arms).
The receiver parameters, trigger threshold, context settings other than the removed
component, and event denominators remain fixed. Every full-arm packet stream is
byte-for-byte identical to its original pilot counterpart. Removing background
reservation releases its quota to trigger capture; this is a practical configuration
ablation under a shared cap, not an equal-expenditure or factorial interaction study.

\begin{table}[htbp]
\centering
\small
% Generated paired development ablations; caps are not equal expenditures.
\begin{tabular}{rlrrr}
\toprule
Cap & Variant & MiB & PCA events & Pre-context (\%) \\
\midrule
256 & Full & 2.63 & 10/11 & 24.2 \\
256 & No pre & 2.63 & 10/11 & 0.0 \\
256 & No post & 2.63 & 10/11 & 24.2 \\
256 & No background & 0.84 & 10/11 & 41.5 \\
512 & Full & 3.57 & 10/11 & 64.3 \\
512 & No pre & 3.55 & 10/11 & 1.8 \\
512 & No post & 3.56 & 10/11 & 64.3 \\
512 & No background & 0.88 & 10/11 & 64.3 \\
1024 & Full & 6.30 & 11/11 & 64.3 \\
1024 & No pre & 6.26 & 11/11 & 1.8 \\
1024 & No post & 6.28 & 11/11 & 64.3 \\
1024 & No background & 0.90 & 10/11 & 64.3 \\
2048 & Full & 11.69 & 11/11 & 67.1 \\
2048 & No pre & 11.66 & 11/11 & 9.1 \\
2048 & No post & 11.68 & 11/11 & 67.1 \\
2048 & No background & 0.90 & 10/11 & 64.3 \\
\bottomrule
\end{tabular}

\caption{Paired development ablations. Pre-context is the macro fraction of original
observations retained in each eligible event's five-minute pre-onset window.
All variants share 11 eligible events; actual expenditure differs.}
\label{tab:ablations}
\end{table}

At the 1,024-byte cap, removing pre-trigger capture preserves PCA evidence for
11/11 events while reducing pre-context coverage from 64.3\% to 1.8\%.
Thus event detection alone conceals a substantial loss of investigation context.
Removing background sampling reduces PCA evidence to 10/11 events and actual
expenditure from 6.30 to 0.90 MiB. Event \texttt{id\_138}, whose annotated span is
53,580 seconds, loses PCA evidence without background sampling. The same event is
lost at the 2,048-byte cap. At 256 bytes, removing background reservation instead
increases pre-context coverage from 24.2\% to 41.5\%, consistent with releasing
capacity for trigger context; the design has budget-dependent tradeoffs.

Removing post-trigger capture leaves PCA event counts and aggregate pre-context
coverage unchanged at all four caps, with small byte savings. This result does not
establish diagnostic value for post-trigger capture, nor its dispensability:
post-event reconstruction and operator diagnosis are not measured. These single
removals do not estimate interactions or establish a globally optimal allocation.
All results remain conditional on the same small development archive.

\paragraph{Independent mission protocol.}
A mission-neutral numerical bundle and a two-stage freeze/run interface now support
an independent evaluation. Before replay, the lock binds input-manifest and code
hashes to the fixed Mission1 retention settings, four caps, and ten seeds. The
new mission's receiver normalization, PCA basis, and score thresholds are calibrated
only on its designated training observations. This evaluates transfer of retention
settings with training-calibrated receivers, not zero-shot receiver transfer.
Training and evaluation annotations are stored separately; evaluation annotations
are parsed only after all packet streams have been captured. The original retention-only interface reports 108 policy,
ablation, and seed arms, together with a full-archive reference. The extended
receiver-mismatch interface is described in Section~\ref{sec:mismatch}.
Descriptive event strata are fixed at durations of at most 60 seconds, over 60
through 600 seconds, and over 600 seconds, and adjacent eligible-observation gaps at a 150-second
boundary (annotation exclusions can contribute to these gaps). Small strata are not independent estimates of generalization uncertainty.

The interface currently requires six synchronized native channels and retains the
existing byte schema. Asynchronous observations and arbitrary channel counts remain
unsupported. Channel selection and split decisions must be documented before
examining evaluation outcomes. Hash locks expose accidental changes but cannot
prove prior label blindness or replace a public experimental record. The interface has been exercised on synthetic fixtures only; an independent-mission result remains pending.
